\documentclass[12pt]{article}
\usepackage[utf8]{inputenc}
\usepackage[T1]{fontenc}
\usepackage{hyperref}
\title{Saudi Arabian Sign Language Language Resources for K-12 Curriculum Development}
\author{Compiled by Claude Sonnet 5.5 (claude-sonnet-5-5) for Andr\'as Kornai}
\date{}
\begin{document}
\maketitle

\section{Introduction}
Saudi Arabian Sign Language (ISO 639-3 \texttt{sdl}; Glottocode saud1238) is the sign language of the Deaf community in Saudi Arabia. The searches found sources describing it as having its own grammar, syntax and vocabulary. Material for it in school settings is limited, and what exists is mostly dictionaries, university courses and research articles.

\section{Primary Educational Resources}
No complete K-12 curriculum in Saudi Arabian Sign Language was found. One report says the dictionary team was also developing a teaching curriculum (Gulf News). Teaching the sign language to hearing students was approved as an optional university course in 2019 at Prince Sattam bin Abdulaziz University, according to the Frontiers article below.
\begin{itemize}
\item Frontiers in Education (2024) article on Saudi Sign Language in education. \url{https://www.frontiersin.org/articles/10.3389/feduc.2024.1245983}
\item Gulf News report on a Deaf Saudi with a master's degree in sign language teaching. \url{https://gulfnews.com/world/gulf/saudi/watch-how-saudi-defied-odds-to-become-first-deaf-person-in-the-region-with-masters-degree-in-sign-language-teaching-1.96973216}
\end{itemize}

\section{Linguistic and Reference Materials}
\begin{itemize}
\item Wikipedia, Saudi Sign Language, which also describes the Saudi Arabia Sign Language Dictionary (2013, Princess Alanood Foundation with Alomary). \url{https://en.wikipedia.org/wiki/Saudi_Sign_Language}
\item King Khalid Foundation page on the ``Integrating the Deaf in Saudi Society'' initiative. \url{https://kka.kkf.org.sa/en/PreviousWinners/Pages/PartnersBranch2013-1.aspx}
\item CSIS analysis of the diverse Arabic sign languages. \url{https://www.csis.org/analysis/reading-signs-diverse-arabic-sign-languages}
\item King Saud University course page referring to the unified Arabic sign dictionary for the Deaf. \url{https://faculty.ksu.edu.sa/ar/salnaji/course/116578}
\end{itemize}

\section{Religious and Community Resources}
\begin{itemize}
\item Saudi Sign Language Association (in Arabic), which works to unify sign language in the Kingdom and licenses interpreters. \url{https://saudipedia.com/%D8%AC%D9%85%D8%B9%D9%8A%D8%A9-%D9%84%D8%BA%D8%A9-%D8%A7%D9%84%D8%A5%D8%B4%D8%A7%D8%B1%D8%A9-%D8%A7%D9%84%D8%B3%D8%B9%D9%88%D8%AF%D9%8A%D8%A9}
\item Deaf Unity page found in search. \url{https://deafunity.org/?p=3404}
\end{itemize}

\section{Digital and Audio Resources}
The Saudi Arabia Sign Language Dictionary is reported to be free online (see the Wikipedia article). No separate video corpus or app was confirmed in the searches.

\section{Curriculum Development Context}
Deaf education in Saudi Arabia is documented mostly in research articles. Work toward a Saudi Arabian Sign Language curriculum would need Deaf-led partners such as the Saudi Sign Language Association and the dictionary authors. Care is needed because some sources treat Saudi Arabian Sign Language and the unified Arabic sign dictionary as separate things, which the sources do not resolve.

\end{document}